% Part: normal-modal-logic
% Chapter: frame-correspondence
% Section: frames

\documentclass[../../../include/open-logic-section]{subfiles}

\begin{document}

\olfileid{nml}{frd}{fra}

\olsection{చట్రాలు}

\begin{defn}
  \emph{చట్రం} అంటే $\mModel{F} = \tuple{W,R}$ అనే జత;
  ఇక్కడ $W$ అనేది లోకాల ఖాళీ కాని సమితి, $R$
  అనేది~$W$పై ద్విస్థానిక సంబంధం. ఏదో మూల్యనిర్ణయం
  $V$కు $\mModel{M} = \tuple{W, R, V}$ అయితేనే,
  నమూనా $\mModel{M}$ చట్రం $\mModel{F} = \tuple{W,R}$పై
  \emph{ఆధారపడుతుంది} అంటాం.
\end{defn}

\begin{defn}
  $\mModel{F}$ చట్రం అయితే, $\mModel{F}$పై ఆధారపడే ప్రతి
  నమూనా~$\mModel{M}$కు $\mSat{M}{!A}$ అయినప్పుడే,
  $!A$ \emph{$\mModel{F}$లో చెల్లుబాటు అవుతుంది,}
  అంటే $\mModel{F} \Entails !A$ అంటాం.

  $\mClass{F}$ చట్రాల వర్గమైతే, ప్రతి చట్రం~$\mModel{F} \in \mClass{F}$కు
  $\mModel{F} \Entails !A$ అయినప్పుడే $!A$
  \emph{$\mClass{F}$లో చెల్లుబాటు అవుతుంది,}
  అంటే $\mClass{F} \Entails !A$ అంటాం.
\end{defn}

చట్రాలు ఆసక్తికరంగా ఉండటానికి కారణం: పథకాలకీ,
ప్రాప్యత సంబంధం~$R$ ధర్మాలకీ మధ్య అనురూపత
\emph{నమూనాల స్థాయిలో కాక, చట్రాల స్థాయిలో} ఉంటుంది.
ఉదాహరణకు, \Ax{T} అన్ని స్వావర్తన నమూనాల్లో సత్యమే;
అయితే \Ax{T} సత్యమైన ప్రతి నమూనా స్వావర్తనమైనది కాదు.
కానీ \Ax{T} అన్ని స్వావర్తన \emph{చట్రాల్లో}
\emph{చెల్లుబాటు} అవడమే కాకుండా, \Ax{T}
చెల్లుబాటు అయ్యే ప్రతి చట్రం స్వావర్తనమైనదన్నది
\emph{నిజం}.

\begin{rem}
చట్రాల వర్గంలో చెల్లుబాటు అనేది నమూనాల వర్గంలో
చెల్లుబాటు అనే భావానికి ఒక ప్రత్యేక సందర్భం:
$\mClass{F} \Entails !A$ అప్పుడే, అప్పుడే
$\mClass{C} \Entails !A$; ఇక్కడ $\mClass{C}$
అనేది~$\mClass{F}$లోని ఏదో చట్రంపై ఆధారపడే
అన్ని నమూనాల వర్గం.

స్పష్టంగానే, \tetoken{సూత్రం}{!!a{formula}} లేదా పథకం
చెల్లుబాటు అయితే, అంటే \emph{అన్ని} నమూనాల వర్గంలో
చెల్లుబాటు అయితే, ఏ చట్రాల వర్గం~$\mClass{F}$లోనైనా
అది చెల్లుబాటు అవుతుంది.
\end{rem}

\end{document}
